% ============================================================
% 第 1 章 引言 —— 赛项背景、任务分解与总体方案
% ============================================================

\section{引言}

\subsection{赛项背景与任务描述}

低空经济已被列为国家战略性新兴产业，智能低空消防既是应急安全的现实刚需，也是低空经济产业化的关键抓手。超高层楼宇、大型商业综合体、地下管廊与化工园区等高危复杂场景不断增多，传统消防装备普遍存在抵达慢、侦察滞后、人员暴露风险高的问题。以四旋翼无人机为载体，融合自主导航、视觉识别、物资投送与态势回传，使其在人员难以抵达的高危区域完成侦察、定位与应急物资投放，是补齐现代消防短板的有效途径。

本赛项要求参赛队自主设计并制作一台智能消防四旋翼无人机，在 $4000\times4000\times3000\unit{mm}$ 的模拟消防场景内，在不使用超宽带定位（UWB）、动作捕捉等外置辅助定位装置的条件下，完成自主起飞、避障穿行、危险品目标识别与标注、消防物资投放、受伤人员识别与标注、救助物资投放以及自主返航降落等任务。场地由起飞区、固定与随机障碍、危险品任务区、人员救助任务区构成，其中随机障碍的位置、危险品警告标志的种类与位置、人员目标的摆放位置均在比赛现场抽签确定，系统必须在没有先验地图与先验目标位置的条件下在线完成感知、决策与飞行。

完整的任务流程可划分为表~\ref{tab:mission_stages} 所示的六个阶段。除起飞与降落外，各阶段之间不存在人工干预，全部由机载系统自主衔接；现场仅允许通过“笔记本电脑一键式”方式启动，起飞后参赛人员不得再触碰计算机，亦不得以任何其他方式遥控飞行器。

\begin{table}[H]
\centering
\caption{智能低空消防任务流程分解}
\label{tab:mission_stages}
\small
\begin{tabular}{L{1.0cm}L{5.0cm}L{6.6cm}}
\toprule
\textbf{阶段} & \textbf{飞行动作} & \textbf{任务与判据} \\
\midrule
1 & 从 $\Phi500\unit{mm}$ 圆形起飞区自主起飞 & 垂直爬升至任务高度，姿态与位置稳定，进入自主模式 \\
2 & 避障飞行至危险品任务区 & 绕过固定障碍与抽签确定的随机障碍，通行间隙不小于 $1300\unit{mm}$ \\
3 & 悬停于危险品任务区上方 & 识别两个任务区内的警告标志，区分正确危险品与干扰项，在回传画面中以红色框标注正确目标 \\
4 & 投放消防物资 & 对准正确目标后投放，物资落点应位于 $400\times400\unit{mm}$ 任务区内 \\
5 & 避障飞行至人员救助任务区并悬停识别 & 识别受伤人员头部目标并以蓝色框标注；决赛第二阶段还需投放救助物资 \\
6 & 自主返航降落 & 返回起飞点降落，螺旋桨停转后计时结束 \\
\bottomrule
\end{tabular}
\end{table}

\subsection{任务分解与关键技术挑战}

将上述任务拆解到技术层面，系统需要同时解决定位、感知、规划、控制与任务决策五类问题。结合场地条件与赛项规则，本文将其归纳为以下五项关键挑战。

\textbf{（1）无外置定位条件下的自主定位与建图。}赛项禁止使用 UWB、动作捕捉等外置辅助定位装置，场地被安全网包围，卫星导航不可用且无先验地图。飞行器必须在飞行中同时完成六自由度位姿估计与障碍物建图，且定位输出需具备足够带宽以支撑 $1\unit{m/s}$ 量级的机动。

\textbf{（2）随机障碍下的在线避障。}固定障碍共 4 处，其中两处宽 $1700\unit{mm}$、高 $2000\unit{mm}$，另两处宽 $500\unit{mm}$、高 $2000\unit{mm}$；随机障碍为两根 $\Phi200\times2000\unit{mm}$ 圆柱，位置在赛前抽签确定，规则仅保证其中一侧通行距离不小于 $1300\unit{mm}$。这意味着系统既不能依赖预置地图，也不能依赖固定航线，必须在飞行中实时感知 $200\unit{mm}$ 直径的细长圆柱并在狭窄通道内完成轨迹重规划。

\textbf{（3）小目标识别与稳定悬停。}危险品任务区与人员救助任务区的平面尺寸仅为 $400\times400\unit{mm}$，任务区内的警告标志与人员头部图像尺寸在数十毫米量级。飞行器悬停高度约 $1.5$--$2.0\unit{m}$ 时，目标在图像中占比小、俯视透视形变大、场地地面为哑光白色导致对比度偏低，且危险品任务区中存在刻意设置的干扰标志，要求识别模型具备小目标召回能力与细粒度分类能力。同时，识别过程要求飞行器保持稳定悬停，使目标始终处于相机视场中心。

\textbf{（4）精准投放。}消防物资与救援物资的单件最大尺寸在 $80\unit{mm}$ 立方范围内，投放落点需落在 $400\times400\unit{mm}$ 的任务区内。投放精度取决于悬停定位精度、目标像素定位精度与投放时机的配合；投放瞬间的质心变化还会对姿态产生扰动，控制器需具备抗扰动能力。

\textbf{（5）系统集成与现场可靠性。}比赛总时长 10 分钟（含调试与两次运行机会），全程仅使用同一块机载电池，且不得更换任何电子元器件；系统需支持自主飞行与手动操控的模式切换，并在自主任务出现故障时由手动模式安全降落；飞行状态需实时回传，运行过程需完整录屏。这些约束对启动流程、故障处理与功耗预算提出了工程化要求。

\subsection{总体技术方案}

针对上述挑战，本文设计并实现了一套“感知—估计—规划—控制—任务”五层闭环的自主飞行系统。硬件平台为 X 型四旋翼，轴距 $450\unit{mm}$，动力采用 F100 KV1100 电机与银燕 45A 电调，搭载 Livox Mid-360 激光雷达、下视广角相机、Jetson Orin NX 16GB 机载计算机与 Pixhawk 6C mini 开源飞控（PX4 固件），并集成双仓舵机拨杆式投放机构，分别装载消防物资与救援物资。

软件层面，系统以 ROS Noetic 为中间件，由四个核心算法模块与一个任务调度模块构成：

\begin{itemize}
    \item \textbf{视觉识别}采用 YOLO26~\cite{jocher2026yolo26}，以端到端无 NMS 的推理方式在机载 GPU 上完成危险品警告标志与人员头部的检测与分类，并通过像素—世界坐标投影实现目标定位与红/蓝框标注；
    \item \textbf{状态估计}采用 FAST-LIO2~\cite{xu2022fastlio2} 激光-惯性紧耦合里程计，融合 Mid-360 点云与飞控 IMU，输出 $200\unit{Hz}$ 高频位姿与世界系点云；
    \item \textbf{规划与避障}采用 Diff-Planner，在概率占据栅格上以 A* 搜索引导路径、以 MINCO 多项式轨迹优化生成满足动力学约束的平滑轨迹，并支持 $10\unit{Hz}$ 在线重规划；
    \item \textbf{飞行控制}采用 PX4Ctrl 几何控制器，基于四旋翼微分平坦特性在 $SO(3)$ 上直接构造期望姿态，配合递推最小二乘油门在线辨识，以 $100\unit{Hz}$ 输出姿态与推力指令至 PX4 飞控~\cite{meier2015px4}；
    \item \textbf{任务调度}由有限状态机实现，负责起飞、巡航、悬停识别、对准投放、返航降落等阶段的切换，并统一处理图传标注、录屏与故障保护。
\end{itemize}

上述模块之间的接口、选型依据、资源预算与工程约束将在第~\ref{sec:tech_stack}~章统一给出。

\subsection{技术报告结构}

本报告围绕智能低空消防任务组织内容。第~\ref{sec:tech_stack}~章给出系统技术栈，包括赛项约束向技术指标的映射、分层架构与数据流、核心模块选型依据、资源预算与备选方案对比。后续章节按视觉感知、状态估计、轨迹规划、飞行控制、任务流程与实机试验的顺序展开详细设计与验证。
